\section{Hardness of calibrating the least penalty}\label{sec:calibration}
This section shows that a sufficient coefficient can be trivial to write
down while a nearly least one is hard to compute. On a binary box with one
relaxed equality, a linear objective, a known unique optimizer and the
known sufficient coefficient $\rho=1$, no polynomial-time algorithm
approximates the least penalty within any polynomial factor unless
$\mathrm P=\mathrm{NP}$; the same holds for the zero-multiplier threshold.
The hardness needs a growing number of binary variables; fixed dimension is
not covered.

Throughout, $m=1$ and the norm is the absolute value. On $\R$ every norm is
a positive multiple of $|\cdot|$, which only rescales the thresholds
(Remark~\ref{rem:rescaling}).

\subsection{The calibration problem}
\begin{definition}[Calibration]\label{def:calibration}
Fix a family of instances of \eqref{eq:model} with rational data, encoded
in binary as in Section~\ref{sec:geometry}, and a function $\varphi$ from the
positive integers to $[1,\infty)$. Given an instance of encoding length $N$
whose least penalty $\rho_*$ is finite, \emph{sufficient calibration within
factor $\varphi$} asks for a rational $\widehat\rho$ with
\begin{equation}\label{eq:onesided}
 \rho_*\le\widehat\rho\le \varphi(N)\,\rho_*,
\end{equation}
and \emph{two-sided calibration within factor $\varphi$} asks for a rational
$\widehat\rho$ with
\begin{equation}\label{eq:twosided}
 \rho_*/\varphi(N)\le\widehat\rho\le \varphi(N)\,\rho_*.
\end{equation}
Calibration at multiplier zero is defined in the same way, with
$\rho_*(0)$ in place of $\rho_*$. \emph{Polynomial-factor calibration}
means that $\varphi=p$ is a fixed polynomial with $p(N)\ge1$ for all $N$.
\end{definition}

A one-sided estimate \eqref{eq:onesided} is itself a sufficient
coefficient: the value-exact coefficients form the interval
$[\rho_*,\infty)$, and the coefficients exact at multiplier zero form
$[\rho_*(0),\infty)$ (Corollary~\ref{cor:ratio} and the discussion after
it). A two-sided estimate need not be sufficient. On the families below a
sufficient coefficient is known in advance, so the only question is how
close to the least one an algorithm can get.

\subsection{A binary-box family}\label{sec:calibration-family}
Take a \textsc{Subset Sum} instance: positive integers $w_1,\ldots,w_t$ and
a positive integer target $\beta$, with encoding length $N_0$. For
$I\subseteq\{1,\ldots,t\}$ write $w(I)=\sum_{i\in I}w_i$. The instance is a
YES instance if $w(I)=\beta$ for some $I$, and a NO instance otherwise.
Deciding this is NP-complete. Append the item $w_{t+1}=\beta+1$, extend the
notation $w(I)$ to $I\subseteq\{1,\ldots,t+1\}$, choose an integer
$\kappa\ge2$, and set
\begin{equation}\label{eq:subsetmodel}
 X=\{0,1\}^{t+2}\ni(x,q),\qquad f(x,q)=-q,\qquad
 r(x,q)=\kappa\sum_{i=1}^{t+1}w_ix_i-(\kappa\beta+1)q.
\end{equation}
The instance is given by the integer coefficient vector
$(\kappa w_1,\ldots,\kappa w_{t+1},-(\kappa\beta+1))$ of $r$; the box and the
objective are implicit.

\paragraph{How the gadget works.}
Read $x$ as the indicator vector of a subset $I$ and $q$ as a switch
whose slice $q=1$ is improving: it lowers the objective from $0$ to $-1$. The residual is
$\kappa w(I)\ge0$ at $q=0$ and $\kappa(w(I)-\beta)-1$ at $q=1$. By
Proposition~\ref{prop:scalar}, the least penalty is governed by how close
the improving slice $q=1$ comes to $r=0$ from each side. Each part of the
construction has one purpose.
\begin{itemize}
 \item The $+1$ in $\kappa\beta+1$ makes the whole slice $q=1$
 equality-infeasible: $r=0$ at $q=1$ would require $\kappa$ to divide
 $\kappa\beta+1$. At $q=0$, $r=0$ forces $x=0$ because all items are
 positive. Hence $(x,q)=(0,0)$ is the unique feasible point and $v=0$,
 known without solving \textsc{Subset Sum}.
 \item The appended item $\beta+1$ fixes the nearest positive residual at
 $q=1$ to $\kappa-1$ on every instance, so the positive side carries no
 information about the instance.
 \item The factor $\kappa$ separates the cases on the negative side: a
 subset with $w(I)=\beta$ gives residual $-1$, while every subset with
 $w(I)\le\beta-1$ gives residual at most $-(\kappa+1)$.
\end{itemize}
Linear optimization over $X$ is trivial; this does not make the penalized
subproblem with the absolute residual easy.

Formally, let $\beta_-=\max\{w(I):I\subseteq\{1,\ldots,t+1\},\
w(I)\le\beta\}$ be the largest attainable sum not exceeding the target
(the empty set shows that it exists), and define
\begin{equation}\label{eq:nearestsums}
 g_-=\kappa(\beta-\beta_-)+1,\qquad g_+=\kappa-1.
\end{equation}
At $q=1$ the residual $\kappa(w(I)-\beta)-1$ is negative exactly when
$w(I)\le\beta$; its magnitude is then $\kappa(\beta-w(I))+1\ge g_-$, with
equality when $w(I)=\beta_-$. The residual is positive exactly when
$w(I)\ge\beta+1$; it is then at least $\kappa-1=g_+$, with equality for
$I=\{t+1\}$. No subset with sum at most $\beta$ contains the item $t+1$, so
$\beta_-$ is the best sum of the original instance. Consequently
\begin{equation}\label{eq:gapcases}
 g_-=1\ \text{in YES instances},\qquad g_-\ge\kappa+1\ \text{in NO instances}.
\end{equation}

\begin{proposition}\label{prop:calibrationdual}
Let $w_1,\ldots,w_t,\beta$ be positive integers, $w_{t+1}=\beta+1$, and
$\kappa\ge2$ an integer. The instance \eqref{eq:subsetmodel} has $v=0$,
attained only at $(x,q)=(0,0)$. For every $\rho\ge0$,
\begin{equation}\label{eq:calibrationdual}
 D_\rho=\min\left\{0,-1+\rho\,\frac{2g_-g_+}{g_-+g_+}\right\},
\end{equation}
and the multiplier $\lambda_\rho=\rho(g_--g_+)/(g_-+g_+)$ attains this
value. Consequently
\begin{equation}\label{eq:calibrationthresholds}
 \rho_*=\frac12\left(\frac1{g_-}+\frac1{g_+}\right),\qquad
 \rho_*(0)=\max\left\{\frac1{g_-},\frac1{g_+}\right\}.
\end{equation}
\end{proposition}
\begin{proof}
The plan is to bound $D_\rho$ above by a two-point balanced mixture and to
show that $\lambda_\rho$ attains this bound. The value $v=0$ and the
uniqueness of the feasible point were shown above.

Let $P_-$ be a point with $q=1$ and $w(I)=\beta_-$, and let $P_+$ be the
point with $q=1$ and $I=\{t+1\}$. Then $(f,r)=(-1,-g_-)$ at $P_-$ and
$(f,r)=(-1,g_+)$ at $P_+$. The weights $g_+/(g_-+g_+)$ on $P_-$ and
$g_-/(g_-+g_+)$ on $P_+$ give a balanced mixture: its average residual is
zero. Its average of $f+\rho|r|$ is $-1+2\rho g_-g_+/(g_-+g_+)$, so
Proposition~\ref{prop:mixture} bounds $D_\rho$ by this number. Weak duality
bounds $D_\rho$ by $v=0$. This proves ``$\le$'' in
\eqref{eq:calibrationdual}.

For the reverse inequality, note that
$\rho-\lambda_\rho=2\rho g_+/(g_-+g_+)\ge0$ and
$\rho+\lambda_\rho=2\rho g_-/(g_-+g_+)\ge0$. A native point with $q=1$ and
$r<0$ has $|r|\ge g_-$, so its penalized value at $\lambda_\rho$ is
\[
 -1+(\rho-\lambda_\rho)|r|\ge-1+(\rho-\lambda_\rho)g_-
 =-1+\rho\,\frac{2g_-g_+}{g_-+g_+}.
\]
A native point with $q=1$ and $r>0$ has $r\ge g_+$, and
$-1+(\rho+\lambda_\rho)r\ge-1+(\rho+\lambda_\rho)g_+$ is the same number.
A native point with $q=0$ has $f=0$ and $r\ge0$, so its value
$(\rho+\lambda_\rho)r$ is nonnegative. Hence $L_\rho(\lambda_\rho)$ is at
least the right-hand side of \eqref{eq:calibrationdual}, and the upper
bound is attained.

By \eqref{eq:calibrationdual}, $D_\rho=0$ exactly when
$\rho\ge(g_-+g_+)/(2g_-g_+)$, which is the formula for $\rho_*$. For
$\rho_*(0)$, apply \eqref{eq:fixed} with $\lambda=0$ and $v=0$: an
infeasible point with $q=0$ contributes $0/|r|=0$, and an infeasible point
with $q=1$ contributes $1/|r|$, which is largest at the smallest residual
magnitudes $g_-$ and $g_+$.
\end{proof}

The formulas agree with Proposition~\ref{prop:scalar}, where
$A_+=1/g_+$ (points with $q=0$ contribute $0$) and $A_-=1/g_-$. At
$\rho=\rho_*$ the interval of exact multipliers in \eqref{eq:scalar}
reduces to the single multiplier
$\lambda_{\rho_*}=(1/g_+-1/g_-)/2$, and both witnesses $P_-$ and $P_+$ tie
with the optimum. So $\rho_*$ is exact at $\lambda_{\rho_*}$ but not
minimizer-exact. Every $\rho>\rho_*$ is minimizer-exact at
$\lambda_{\rho_*}$: then $\rho-\lambda_{\rho_*}>1/g_-$ and
$\rho+\lambda_{\rho_*}>1/g_+$, so every infeasible point has a positive
penalized value. On this family $1/(2\kappa)<\rho_*\le1$, because
$g_+=\kappa-1$ and $g_-,g_+\ge1$.

\paragraph{A sufficient coefficient is known.}
For every instance of the family, $\rho=1$ is exact at multiplier zero.
Every nonzero residual is an integer and the objective is $0$ or $-1$, so
$f+|r|\ge0$ at every native point. An infeasible point has $f+|r|=0$ only
if $q=1$ and $|r|=1$. The residual $-1$ occurs exactly in YES instances,
and the residual $+1$ occurs exactly when $\kappa=2$ (through
$I=\{t+1\}$). These ties are the only obstruction: every $\rho>1$ is
minimizer-exact at multiplier zero.

\paragraph{The YES/NO gap.}
By \eqref{eq:gapcases} and \eqref{eq:calibrationthresholds},
\begin{equation}\label{eq:thresholdgap}
 \begin{aligned}
 &\rho_*=\frac{\kappa}{2(\kappa-1)}>\frac12,\qquad
 \rho_*(0)=1 &&\text{in YES instances},\\
 &0<\rho_*\le\frac{\kappa}{\kappa^2-1},\qquad
 \rho_*(0)=\frac1{\kappa-1} &&\text{in NO instances}.
 \end{aligned}
\end{equation}
The ratio between the YES value and the NO bound is $(\kappa+1)/2$ for
$\rho_*$ and $\kappa-1$ for $\rho_*(0)$. For $\kappa=2$ the zero-multiplier
threshold equals $1$ in both cases, so part~(b) below uses $\kappa\ge3$.

\subsection{Testing exactness and calibrating are hard}
The first result shows that testing a single given coefficient is already
hard. The membership argument uses that $v=0$ is known on this family; for
a general instance, certifying $D_\rho<v$ would also require
information about $v$.

\begin{theorem}\label{thm:coNP}
Consider the instances \eqref{eq:subsetmodel} with $\kappa=4$, each given
by the integer coefficient vector $(4w_1,\ldots,4w_{t+1},-(4\beta+1))$ of
$r$. Deciding whether $\rho=1/3$ is value-exact, that is, whether
$D_{1/3}=0$, is coNP-complete. More precisely, $D_{1/3}=-1/2$ for YES
instances of \textup{\textsc{Subset Sum}} and $D_{1/3}=0$ for NO instances.
\end{theorem}
\begin{proof}
For YES instances, $g_-=1$ and $g_+=3$, and \eqref{eq:calibrationdual} gives
$D_{1/3}=\min\{0,-1+\tfrac13\cdot\tfrac64\}=-1/2$. For NO instances,
$\rho_*\le4/15<1/3$ by \eqref{eq:thresholdgap}, so $D_{1/3}=0$.

Hardness: the map from a \textsc{Subset Sum} instance to the coefficient
vector is computable in polynomial time and sends NO instances to exact
coefficients and YES instances to non-exact ones. It is therefore a
polynomial reduction from the complement of \textsc{Subset Sum}, which is
coNP-complete.

Membership in coNP: whether an input vector belongs to the family can be
checked in polynomial time, since $\beta$, the items $w_i$ and the
condition $w_{t+1}=\beta+1$ are recovered by integer division by $4$. On a
member of the family, $D_{1/3}<0$ holds exactly in YES instances, and a
subset $I\subseteq\{1,\ldots,t\}$ with $w(I)=\beta$ is a polynomial
certificate of this.
\end{proof}
Consequently, computing $\rho_*$ on this family is NP-hard, and evaluating
$D_{1/3}$ within absolute error strictly below $1/4$ decides
\textsc{Subset Sum}. In contrast, $D_0=-1$ on every instance: without the
norm term the Lagrangian dual is trivial to evaluate and never closes the
gap.

The second result says that even approximating the threshold, from above,
within a polynomial factor is hard.

\begin{theorem}\label{thm:polyfactor}
Let $p$ be a fixed polynomial with $p(N)\ge1$ for all $N$. Unless
$\mathrm P=\mathrm{NP}$, no polynomial-time algorithm achieves sufficient
calibration within factor $p$ on the family \eqref{eq:subsetmodel}, either
\begin{enumerate}
 \item[(a)] for the least penalty: on every instance of encoding length
 $N$, return a rational $\widehat\rho$ with
 \begin{equation}\label{eq:polyfactor}
  \rho_*\le\widehat\rho\le p(N)\,\rho_*;
 \end{equation}
 \item[(b)] or at multiplier zero: on every instance of encoding length $N$,
 return a rational $\widehat\rho$ with
 $\rho_*(0)\le\widehat\rho\le p(N)\,\rho_*(0)$, where
 $\rho_*(0)=\max\{1/g_-,1/g_+\}$; this equals $1$ in YES instances and
 $1/(\kappa-1)$ in NO instances.
\end{enumerate}
\end{theorem}
\begin{proof}
The plan is to choose $\kappa$ exponentially large, so that the YES/NO
ratio in \eqref{eq:thresholdgap} exceeds $p(N)$, and then to read off the
\textsc{Subset Sum} answer by comparing $\widehat\rho$ with a fixed number.

Given a \textsc{Subset Sum} instance of length $N_0$, take
$\kappa=2^{N_0}$. The constructed instance has $t+2\le N_0+2$ coefficients,
each with $O(N_0)$ bits, so its length is $N=O(N_0^2)$ and it is computed
in polynomial time. In particular $p(N)$ is bounded by a polynomial in
$N_0$.

(a) In YES instances every admissible output satisfies
$\widehat\rho\ge\rho_*>1/2$. In NO instances every admissible output
satisfies
\[
 \widehat\rho\le p(N)\,\frac{\kappa}{\kappa^2-1}
 =p(N)\,\frac{2^{N_0}}{4^{N_0}-1}<\frac12
\]
for all sufficiently large $N_0$, since the exponential term dominates.
(b) For $N_0\ge2$ we have $\kappa\ge3$. In YES instances every admissible
output satisfies $\widehat\rho\ge1$; in NO instances
$\widehat\rho\le p(N)/(2^{N_0}-1)<1$ for all sufficiently large $N_0$.

In both cases there is a constant $n_1$, depending only on $p$, such that
the comparison of $\widehat\rho$ with $1/2$ (respectively $1$) decides
every \textsc{Subset Sum} instance with $N_0\ge n_1$. Instances with
$N_0<n_1$ are finitely many and are decided by exhaustive enumeration. The
output $\widehat\rho$ of a polynomial-time algorithm has polynomially many
bits, so the comparison takes polynomial time. This would decide
\textsc{Subset Sum} in polynomial time.
\end{proof}

\begin{remark}[Two-sided estimates and constant factors]\label{rem:twosided}
Two-sided calibration is not easier. If $\widehat\rho$ satisfies
\eqref{eq:twosided} with $\varphi=p$, replace $p$ by a pointwise larger polynomial
$p'$ with positive integer coefficients (so that $p'(N)$ is an integer
computable in polynomial time); then $p'(N)\widehat\rho$ satisfies
\eqref{eq:onesided} with $\varphi=p'^2$, again a polynomial. Thus
Theorem~\ref{thm:polyfactor} also rules out two-sided polynomial-factor
calibration, in both parts. For a fixed constant factor $C\ge1$, an
exponential $\kappa$ is not needed. In part~(a), a one-sided estimate within
factor $C$ separates the YES value $\kappa/(2(\kappa-1))$ from the NO bound
$C\kappa/(\kappa^2-1)$ exactly when $C<(\kappa+1)/2$, so any constant
$\kappa$ with $\kappa+1>2C$ suffices. In part~(b), any constant $\kappa$ with
$\kappa-1>C$ suffices.
\end{remark}

\begin{remark}[Subexponential factors]\label{rem:subexpfactor}
The proof gives more than polynomial factors. For every fixed
$\epsilon\in(0,1)$, sufficient calibration within factor
$\varphi(N)=2^{N^{1-\epsilon}}$ is also impossible in polynomial time unless
$\mathrm P=\mathrm{NP}$, in both parts. Choose an integer
$c>(1-\epsilon)/\epsilon$ and take $\kappa=2^{N_0^c}$. Each coefficient
then has at most $N_0^c+O(N_0)$ bits, so $N=O(N_0^{c+1})$ and the
construction is polynomial. Since $(c+1)(1-\epsilon)<c$, we have
$N^{1-\epsilon}\le N_0^c-2$, that is $\varphi(N)\le\kappa/4$, for all
sufficiently large $N_0$. Then in part~(a) every NO output satisfies
$\widehat\rho\le \varphi(N)\kappa/(\kappa^2-1)\le\kappa^2/(4(\kappa^2-1))<1/2$,
while every YES output exceeds $1/2$. In part~(b) every NO output satisfies
$\widehat\rho\le \varphi(N)/(\kappa-1)\le\kappa/(4(\kappa-1))<1$, while every YES
output is at least $1$.
\end{remark}

\begin{remark}[Rescaling does not help]\label{rem:rescaling}
Replacing $r$ by $cr$ for a rational $c>0$ replaces $L_\rho(\lambda)$ by
$L_{c\rho}(c\lambda)$ of the original instance; replacing the absolute value
by the norm $c|\cdot|$ replaces it by $L_{c\rho}(\lambda)$. In both cases $\rho_*$ and $\rho_*(0)$ are divided by $c$, and the
YES/NO ratios $(\kappa+1)/2$ and $\kappa-1$ are unchanged. A guarantee
within a factor is invariant under such rescaling, so no choice of scaling
makes calibration easier on this family.
\end{remark}

\begin{remark}[Weak NP-hardness]\label{rem:weak}
The reduction is from \textsc{Subset Sum}, which is NP-hard only when its
numbers are large, and the constructed coefficients are large as well. This
is unavoidable on this family: a dynamic program over the attainable values
of $\sum_iw_ix_i$ computes $g_-$, hence $\rho_*$ and $\rho_*(0)$ exactly, in
time polynomial in $N$ and the largest coefficient magnitude. So the family
gives weak NP-hardness only. Proposition~\ref{prop:stableset} below gives
strong NP-hardness of exact computation with coefficients in
$\{-1,0,1\}$, at the price of a native set over which linear optimization
is itself hard.
\end{remark}

\begin{remark}[Where the hardness lies]\label{rem:hardnesslocation}
The difficulty is testing whether a given coefficient is exact
(Theorem~\ref{thm:coNP}), not locating the threshold once such a test is
available. Suppose an oracle decides, for a given instance of the family and
a rational $\rho$, whether $\rho$ is value-exact. Since
$\rho_*\in(1/(2\kappa),1]$ and the value-exact coefficients form
$[\rho_*,\infty)$, testing $\rho=2^{-j}$ for
$j=0,1,\ldots,\lceil\log_2\kappa\rceil+1$ finds the largest $j$ with
$2^{-j}$ exact; then $2^{-j}\ge\rho_*>2^{-j-1}$, so $\widehat\rho=2^{-j}$ is
a sufficient calibration within factor two, after $O(\log\kappa)=O(N)$
oracle calls. The same works at multiplier zero with an oracle for
$L_\rho(0)=v$. In decomposition methods the penalized subproblem is solved
anyway. The results above say that nearly least penalties are as hard to
certify as those subproblems, even when the original problem is trivial.
\end{remark}

\subsection{Unit coefficients with a hard native set}
The binary-box family has large numbers. The next family has all
coefficients in $\{-1,0,1\}$, and its least penalty is the stability number
of a graph. This gives strong NP-hardness of exact computation, because the
native set now encodes a hard optimization problem.

For a simple graph $G=(V,E)$ with $V\ne\varnothing$, let $\alpha(G)$ be the
maximum size of a stable set. Take binary variables $x_i$ ($i\in V$),
$y_+$ and $y_-$, with native constraints
\begin{equation}\label{eq:stablemodel}
 y_++y_-\le1,\qquad x_i\le y_++y_-\ (i\in V),\qquad
 x_i+x_j\le1\ (ij\in E),
\end{equation}
objective $f=-\sum_{i\in V}x_i$ and relaxed equality $r=y_+-y_-=0$.

\begin{proposition}\label{prop:stableset}
For the instance \eqref{eq:stablemodel}, $v=0$ is attained only at the
zero vector, and for all $\rho\ge0$ and $\lambda\in\R$,
\begin{equation}\label{eq:stabledual}
 L_\rho(\lambda)=\min\{0,\rho-\alpha(G)-|\lambda|\},\qquad
 D_\rho=\min\{0,\rho-\alpha(G)\},\qquad \rho_*=\rho_*(0)=\alpha(G).
\end{equation}
At $\rho=\alpha(G)$, $\lambda=0$ is the only exact multiplier and the
infeasible points given by maximum stable sets tie with the optimum. At
multiplier zero, $\rho$ is minimizer-exact if and only if $\rho>\alpha(G)$.
Consequently, computing $\rho_*$ is strongly NP-hard. Given $G$ and an
integer $c$ as input, deciding whether $\rho=c$ is value-exact, that is,
whether $\alpha(G)\le c$, is coNP-complete; a stable set of size $c+1$
certifies failure. For each fixed $c$, enumerating the $(c+1)$-subsets of
$V$ decides it in polynomial time.
\end{proposition}
\begin{proof}
If $r=0$, then $y_+=y_-$, and $y_++y_-\le1$ forces $y_+=y_-=0$; then
$x_i\le0$ for all $i\in V$. So the zero vector is the unique feasible point and
$v=0$. The native points form three slices. If $y_+=y_-=0$, then $x=0$
and $(f,r)=(0,0)$. If $(y_+,y_-)=(1,0)$, the constraints on $x$ say exactly
that $x$ is the indicator vector of a stable set $S$, and
$(f,r)=(-|S|,1)$. If $(y_+,y_-)=(0,1)$, likewise $(f,r)=(-|S|,-1)$.
Minimizing $f+\lambda r+\rho|r|$ over the three slices gives
$\min\{0,\rho-\alpha(G)+\lambda,\rho-\alpha(G)-\lambda\}$, which is the
formula for $L_\rho(\lambda)$. Its supremum over $\lambda$ is attained at
$\lambda=0$, which gives $D_\rho$, and $D_\rho=0$ exactly when
$\rho\ge\alpha(G)$; the same holds for $L_\rho(0)$, so
$\rho_*=\rho_*(0)=\alpha(G)$.

At $\rho=\alpha(G)$, $L_\rho(\lambda)=\min\{0,-|\lambda|\}$ equals $0$ only
for $\lambda=0$, and a maximum stable set with $y_+=1$ then has penalized
value $0$. At multiplier zero, every infeasible point has value
$\rho-|S|$ for a stable set $S$, which is positive for all such points
exactly when $\rho>\alpha(G)$.

All coefficients and right-hand sides lie in $\{-1,0,1\}$, and computing
$\alpha(G)$ is NP-hard, so computing $\rho_*$ is strongly NP-hard.
Deciding $\alpha(G)\le c$ for an input integer $c$ is the complement of the
NP-complete stable-set decision problem, and a stable set of size $c+1$ is
a polynomial certificate that $\rho=c$ is not value-exact. For fixed $c$
there are $O(|V|^{c+1})$ subsets of size $c+1$ to check.
\end{proof}
Unlike the binary box, where linear optimization over $X$ is trivial, here
linear optimization over the native set is already the maximum stable-set
problem. We use this family only for strong NP-hardness of exact
computation. It cannot exclude every polynomial factor as
Theorem~\ref{thm:polyfactor} does: since $1\le\alpha(G)\le|V|\le N$, the output $\widehat\rho=|V|$ is
always a sufficient calibration within factor $|V|\le N$. The numeric
family \eqref{eq:subsetmodel} is what yields arbitrary polynomial factors.

\subsection{Relation to prior work}
In their conclusion, \citet{LefebvreSchmidt2025} ask whether the smallest
penalty parameter that closes the duality gap can be computed in
polynomial time, and conjecture a negative answer. Their gap is defined
through the supremum over multipliers (their Definition~1), which is
value-exactness here, and their question concerns MILPs and MIQPs. The
binary-box family \eqref{eq:subsetmodel} is a pure binary integer linear
program, so it lies within the scope of their question.
Theorem~\ref{thm:polyfactor}(a) confirms their conjecture unless
$\mathrm P=\mathrm{NP}$, even for approximation within any polynomial
factor, in growing binary dimension. Section~\ref{sec:related} compares
the results with earlier hardness results for choosing penalties.
